\section{Comments}

Your current discussion mixes several different operations that actually serve different purposes. The key point is that mean removal, spike suppression, and tail lifting (e.g., whitening) are three different operations with three different objectives, and they should not be conflated.

Subtracting the mean only removes the rank-1 DC offset of the representation. This is useful in cases where the mean acts as a shared non-informative carrier—for example in MoE routing or low-bit quantization, where a large common shift wastes dynamic range or dominates linear projections. In such cases, centering improves the signal-to-noise ratio without changing the overall geometry. However, mean removal does not solve the broader anisotropy problem, because the spectrum can still remain highly skewed: a few dominant singular directions may still concentrate most of the energy while the long tail remains weak. Therefore, mean removal should be viewed mainly as a bias correction, not a general solution to representation anisotropy.

Whitening is a much stronger operation. In addition to removing the mean, it rescales the covariance so that different directions have comparable variance, effectively lifting the long-tail directions relative to dominant principal modes. This makes whitening a more natural tool if the goal is to treat anisotropy in a general sense. However, whitening can also distort useful structure, amplify noise, and destroy meaningful non-isotropic organization in the representation space. In many cases the dominant directions may contain useful shared features rather than purely nuisance bias.

Because of this, the “mean issue” is not a single problem with a single solution. The effect of the mean depends on the downstream mechanism. In quantization the mean wastes representational range, so separating it helps. In routing the mean may dominate a linear selector, so centering or spike attenuation may help. But in optimization or representation learning the mean may be partly harmful and partly useful because it contributes both to curvature concentration and to shared feature organization.

Therefore, the right framework should distinguish analysis space from intervention space. Analytically, you can decompose the representation into the mean component, dominant singular directions, and the long tail. But intervention should be task-specific: centering for quantization or routing, low-rank suppression for curvature concentration, and only cautious whitening when full spectral rebalancing is truly necessary.

In other words, mean removal is a local surgical tool, while whitening is a global geometric transform. If the goal is to build a general theory, the object of study should not be the mean alone but the broader spike–tail geometry and how different system components interact with it. 